\subsection*{Other proofs of fundamental theorems.}
\phantomsection
\addcontentsline{toc}{subsection}{Other proofs of fundamental theorems.}

In {[}278 b{]}, F. Schur shows that in canonical co-ordinates the \(\psi_{ik}\) of (15) satisfy the differential equations

\[
\frac {d}{d t} (t \psi_ {i k} (t a)) = \delta_ {i k} + \sum_ {j, l} c _ {j l} ^ {k} t a _ {l} \psi_ {j i} (t a).\tag{25.23}
\]

These integrate and give a formula equivalent to the formula

\[
\varpi (X) = \sum_ {n \geq 0} \frac {1}{(n + 1) !} (\operatorname{ad} (X)) ^ {n}\tag{25.24}
\]

for the right differential \(\varpi(X)\) of the exponential map at the point X; in particular, in canonical co-ordinates, the \(\psi_{ij}\) are extended into entire functions of the \(a_{k}\) . F. Schur deduces from this a result making precise an earlier remark of Lie: if, in the definition (25.4) of groups of transformations, it is only assumed that the \(f_{i}\) are in the class \(C^{2}\) , then the group is holoedrically isomorphic to an analytic group. \(^{12}\) Following his research on the integration of differential systems, E. Cartan ({[}52 a{]}, v. II₂, p.~371) introduces in 1904 the Pfaffian forms

\[
\omega_ {k} = \sum_ {i = 1} ^ {r} \psi_ {k i} d a _ {i} (1 \leq i \leq r)\tag{25.25}
\]

(with the notation of (25.15)), later called Maurer-Cartan forms. The conditions (25.17) of Maurer can then be written

\[
d \omega_ {k} = - \frac {1}{2} \sum_ {i, j} c _ {i j} ^ {k} \omega_ {i} \wedge \omega_ {j};
\]

E. Cartan shows that it is possible to develop the theory of finite continuous groups starting with the \(\omega_{k}\) and establishes the equivalence of this point of view and that of Lie. But, for him, the interest in this method is especially that it adapts itself to ``infinite and continuous'' groups for which he pushes the theory much further that had been done by Lie, and that it allows the setting up of his theory of the generalised ``mobile marker''.

